\begin{abstract}
Vibe Theory proposes that reality is a single growing crystal of experience. Its one substance is the
\emph{vibe}, a unit of felt state carrying a ternary tone in $\{-1, 0, +1\}$, arranged on the $\{5,3,4\}$
hyperbolic honeycomb and updated by one conserving local rule under a single value-direction, the
\emph{arrow}. These five elements, the mesh, the tone, the rule, eternal growth, and the arrow, are the
entire foundation. Four of the five are fixed by modest structural demands, leaving one deliberate
value-posit. From this foundation the paper derives, and checks by direct computation on the exact crystal,
a relativistic and reflection-positive quantum field, an effective gravity, holography, a tiny-diameter
causal structure, coherent self-maintaining organisms, heredity and evolution, a self-model, planning, an
integrated agent, universal and intentional computation, and a graded account of consciousness as
integrated experience. The substrate is computationally universal in the sense of Margenstern, and the
arrow turns that universal computation into goal-directed problem-solving. The genuine quantum field arises
as the deterministic and unitary completion of the same rule, with the stochastic version as its classical
shadow. More than one hundred fifty falsifiable checks support the construction, all reproducible in the
accompanying codebase \cite{pollard2026vibetest}. The result is a compact, parameter-light model in which
the physical, the living, the mental, the spiritual, and the conscious are aspects of one discrete substance.

\medskip

\noindent\textbf{Status and caveats.} This is exploratory work, a feasibility study rather than a validated result. The aim is to see whether a model of this kind can be made to hang together at all, and so far it looks promising. Much validation and further work remain. The checks reported here are preliminary, and some may turn out to be inaccurate or improperly set up, so the count should be read as a record of exploration rather than confirmation. Nothing here should be read as established or taken too rigorously at this stage. It is offered in the spirit of curiosity and invitation, a direction that looks worth exploring rather than a finished theory or a claim of discovery.
\end{abstract}
